\def\EditionMode{1}
\def\EditionLabel{Local math-prose proof}
\def\EditionStatus{Proposed wrappers; canonical source unchanged}
\documentclass[10pt,oneside]{book}
\usepackage{fontspec}
\setmainfont{Latin Modern Roman}
\setsansfont{Latin Modern Sans}
\usepackage{amsmath,amssymb,mathtools}
\usepackage{unicode-math}
\setmathfont{Latin Modern Math}
\usepackage{xcolor,tikz-cd,lua-ul,marginnote,pdflscape}
\usetikzlibrary{arrows.meta,decorations.pathmorphing,calc,shapes.geometric}
\tikzcdset{ampersand replacement=\&}
\tikzcdset{every label/.append style={text=black}}
\usepackage[a4paper,left=18mm,right=45mm,top=18mm,bottom=22mm,marginparwidth=28mm,marginparsep=7mm]{geometry}
\usepackage[unicode,colorlinks=true,urlcolor=black,linkcolor=black]{hyperref}
\usepackage{bookmark}
\setlength{\parindent}{0pt}
\setlength{\parskip}{4pt plus 1pt minus 1pt}
\setlength{\emergencystretch}{1.5em}
\raggedbottom
% Include final-line descenders inside the declared text-height rectangle.
\maxdepth=0pt
\definecolor{PSHand}{RGB}{25,58,132}
\definecolor{PSDelete}{RGB}{146,37,37}
\definecolor{PSTypedOld}{RGB}{95,95,95}
\ifnum\EditionMode=2\colorlet{blue}{black}\else\colorlet{blue}{PSHand}\fi
\RenewDocumentCommand{\textcircled}{m}{\tikz[baseline=(PSCircle.base)]\node[draw,circle,inner sep=.5pt,line width=.4pt](PSCircle){#1};}
\newif\ifPSInApparatus
\newcounter{psnote}
\newcounter{psvariant}
\newcounter{pslinkmark}
\newif\ifPSInLink
\newsavebox{\PSNoteBox}
\newdimen\PSNoteLimit
{\footnotesize\global\PSNoteLimit=3\baselineskip}
\newcommand{\CurrentScanURL}{}
\newcommand{\CurrentKey}{}
\newcommand{\CurrentScan}{}
\newcommand{\pssource}[3]{\psflushmargins\clearpage\gdef\CurrentKey{#1}\gdef\CurrentScan{#2}\gdef\CurrentScanURL{#3}\setcounter{psvariant}{0}\typeout{PS-SOURCE: #1}}
\newcommand{\psdivision}[2]{\pdfbookmark[0]{#1}{#2}\par\noindent{\large\bfseries #1\par}\medskip}
\newcommand{\pslinkmarker}[2]{\stepcounter{pslinkmark}\latelua{texio.write_nl('log','PS-LINK-SHIP|#1|\CurrentKey|\arabic{pslinkmark}|#2')}}
\newcommand{\pslink}[2]{\ifPSInLink\pslinkmarker{#1}{nested}#2\else\begingroup\PSInLinktrue\href{\CurrentScanURL}{\pslinkmarker{#1}{outer}#2}\endgroup\fi}
\newcommand{\scanpage}[1]{\ifnum\EditionMode=2\leavevmode\marginpar{\scriptsize\sffamily\pslink{scanpage}{scan #1}}\else\par\noindent{\small\sffamily\pslink{scanpage}{[scan #1]}}\par\medskip\fi}
\newcommand{\mspage}[1]{\ifnum\EditionMode=2\leavevmode\marginpar{\scriptsize\sffamily p.\,#1}\else{\small\sffamily[p.\,#1]}\fi}
\newcommand{\hand}[1]{{\ifnum\EditionMode=2\color{black}\else\color{PSHand}\fi #1}}
\let\over\relax
\NewDocumentCommand{\over}{O{} m m}{{\textcolor{PSTypedOld}{#2}\ensuremath{^{\!\nearrow}}\if\relax\detokenize{#1}\relax\hand{#3}\else #3\fi}}
\NewDocumentCommand{\inserted}{O{} m}{{\if\relax\detokenize{#1}\relax\hand{\ensuremath{\langle}#2\ensuremath{\rangle}}\else\ensuremath{\langle}#2\ensuremath{\rangle}\fi}}
\newcommand{\deleted}[1]{{\color{PSDelete}\ifmmode\text{[del.: }#1\text{]}\else[\textsf{del.:} #1]\fi}}
\newcommand{\xout}[1]{{\color{PSTypedOld}\ifmmode\text{[x: }#1\text{]}\else[\textsf{x:} #1]\fi}}
\let\PSMathUnderline\underline
\RenewDocumentCommand{\underline}{+m}{\ifmmode\PSMathUnderline{#1}\else\underLine{#1}\fi}
\NewDocumentCommand{\uline}{+m}{\underline{#1}}
\NewDocumentCommand{\uuline}{+m}{\ifmmode\PSMathUnderline{\PSMathUnderline{#1}}\else\underLine{\underLine{#1}}\fi}
\newunderlinetype\PSWaveLine{\leaders\hbox{\lower.7ex\hbox{\tikz[x=1pt,y=1pt,baseline=0pt]{\draw[line width=.3pt] (0,0) .. controls (1,1) and (2,1) .. (3,0) .. controls (4,-1) and (5,-1) .. (6,0);}}}}
\NewDocumentCommand{\uwave}{+m}{{\PSWaveLine #1}}
\newcommand{\unsure}[1]{\pslink{unsure}{[#1\ifmmode\,?\else\,?\fi]}}
\newcommand{\illegible}[1]{\ifmmode\text{\pslink{illegible}{[illeg.\ #1]}}\else\pslink{illegible}{[illeg.\ #1]}\fi}
\NewDocumentCommand{\diagram}{+m}{\ifmmode\vcenter{\hbox{#1}}\else\par\medskip{\centering #1\par}\medskip\psflushmargins\fi}
\newenvironment{diagramblock}{\par\begin{center}}{\end{center}\par}
\newcommand{\PSMathNotes}{}
\makeatletter
\NewDocumentCommand{\psplacednote}{m +m}{%
      \sbox{\PSNoteBox}{\parbox{\marginparwidth}{\footnotesize\raggedright\hand{\textsf{m#1.} #2}}}%
      \ifdim\dimexpr\ht\PSNoteBox+\dp\PSNoteBox\relax<\PSNoteLimit
        \leavevmode\marginpar{\usebox{\PSNoteBox}}%
      \else
        \par\nobreak\smallskip\hrule\nobreak\smallskip
        {\footnotesize\textsf{Marginal note m#1.}\quad\hand{#2}\par}
        \nobreak\smallskip\hrule\smallskip
      \fi}
\newcommand{\psflushmargins}{\ifmmode\else\ifinner\else\ifx\PSMathNotes\@empty\else\let\PSTempNotes\PSMathNotes\global\let\PSMathNotes\@empty\PSTempNotes\fi\fi\fi}
\NewDocumentCommand{\margin}{+m}{%
  \ifPSInApparatus\ifmmode\hand{\text{[marg.: }#1\text{]}}\else\hand{\textsf{[marg.: }#1\textsf{]}}\fi
  \else\stepcounter{psnote}%
    \ifmmode
      \text{\textsuperscript{m\arabic{psnote}}}%
      \begingroup\edef\PSQueue{\endgroup\noexpand\g@addto@macro\noexpand\PSMathNotes{\noexpand\psplacednote{\arabic{psnote}}{\unexpanded{$#1$}}}}\PSQueue
    \else
      \textsuperscript{m\arabic{psnote}}\psplacednote{\arabic{psnote}}{#1}%
    \fi
  \fi}
\AtBeginDocument{\everymath\expandafter{\the\everymath\aftergroup\psflushmargins}\everydisplay\expandafter{\the\everydisplay\aftergroup\psflushmargins}}
\makeatother
\newcommand{\psnumber}[2]{\expandafter\gdef\csname psnum@#1\endcsname{#2}}
\NewDocumentCommand{\var}{+m +m m m}{\ifnum\EditionMode=0 #1\else\ifnum\EditionMode=1\ifmmode\text{\textsuperscript{\tiny\csname psnum@#4\endcsname}}\else\textsuperscript{\tiny\csname psnum@#4\endcsname}\fi\fi #2\fi}
\newcommand{\psapparatusstart}{\par\bigskip\hrule\smallskip\PSInApparatustrue\begingroup\footnotesize\raggedright\textsf{Apparatus for scan \CurrentScan}\par}
\newcommand{\psapparatusend}{\par\endgroup\PSInApparatusfalse}
\NewDocumentCommand{\psentry}{m m m +m +m m}{\par\noindent\textbf{#2. #3.}\quad\textit{Typescript:} #4\quad\textit{Corrected:} \if\relax\detokenize{#5}\relax\textnormal{[omitted]}\else #5\fi\quad #6\par}
\newcommand{\psfrontmatter}{%
\hypersetup{pageanchor=false}\begin{titlepage}\raggedright
{\Large Alexander Grothendieck\par}\vspace{12mm}
{\Huge Pursuing Stacks\par}\medskip{\LARGE 1983\par}\vspace{12mm}
{\Large\EditionLabel\par}\vspace{8mm}
{\large\EditionStatus\par}\vfill
Université de Montpellier, Grothendieck archive, cotes 134-2--134-8.\par
Including the 26 source pages of R. Brown's translation at cote 134-2, sites 115--143.\par
\vspace*{2mm}\end{titlepage}\hypersetup{pageanchor=true}
\chapter*{Editorial statement}
AI-transcribed by Claude Opus 5.5 (Anthropic, effort high); checked against the scans, and this edition built, by GPT-6 Astra (OpenAI, ChatGPT Pro mode); prompts and records by Claude (Anthropic).

The witness consulted and not reproduced is the LaTeX edition of Mateo Carmona with the collaboration of Ulrik Buchholtz (arXiv 2111.01000 v2, 2021), extended from the 2015 edition by ``the scrivener'', CC0 1.0.

The translation in the appendix is R. Brown's, prepared with assistance from L. Breen and J-L. Loday, as his note says; the rights in it belong to its translator or his heirs.

The rights in the original work belong to its author or the author's representatives or heirs. Where we hold rights, in the AI-produced transcription, translation, corrections and typesetting, they are dedicated to the public domain under CC0 1.0. This edition is published for educational, non-profit purposes, and we claim no rights we do not hold. If you hold rights in this work and object to its publication here, please open a GitHub issue and we will take it down.

\section*{How to read this form}
\ifnum\EditionMode=0
The diplomatic form preserves typed and handwritten wording, cancellations, insertions, overwritten readings, underlining, doubts, and original page numbers. Blue marks handwriting; red labels handwritten cancellations; gray labels typewriter cancellations and underlying overwritten text. Angle brackets mark insertions; a superscript arrow separates an overwritten reading from its replacement. Typed insertions and typed overwriting remain plain. The disposition of long marginal notes is identified explicitly.
\else
The corrected text applies the author's insertions, replacements, cancellations, and recorded transpositions. It retains his English and notation, including consistent forms such as ``standart'', ``developped'', ``it's'' for ``its'', and ``A la''. An isolated slip is corrected only when specifically recorded and justified. Mathematical anomalies are preserved and discussed in the guide.
\ifnum\EditionMode=1
Blue continues to identify handwriting. A numbered call identifies each maximal editorial operation, with its full diplomatic and corrected readings and reason in the apparatus following the source page. Internal layers of a wholly replaced or deleted passage are preserved within that entry rather than duplicated as separate entries. These entries are generated from the same variants that govern the running text.
\else
This reading form omits the apparatus, cancelled material, superseded readings, insertion brackets, and color distinctions. It retains the resulting text, marginal notes, underlining, and unresolved doubts. Scan and source-page references appear in the margin.
\fi\fi

Every scan mark and remaining doubt links to its source page in the archive viewer. Doubts are editorial reading marks, not emendations. The source-page breaks identify the archival sequence; a source page may occupy more than one printed page. A partial working edition is explicitly labelled and contains only reviewed source pages.
\clearpage}

\begin{document}
\pssource{ps_134-7_054}{134-7/54}{https://grothendieck.umontpellier.fr/web/viewer.html?file=../134-7.pdf\#page=55}
\psnumber{ps_134-7_054-v0001}{1}
\psnumber{ps_134-7_054-v0002}{2}
\psnumber{ps_134-7_054-v0003}{3}
\psnumber{ps_134-7_054-v0004}{4}
\psnumber{ps_134-7_054-v0005}{5}
\psnumber{ps_134-7_054-v0006}{6}
\psnumber{ps_134-7_054-v0007}{7}
\psnumber{ps_134-7_054-v0008}{8}
\psnumber{ps_134-7_054-v0009}{9}
\psnumber{ps_134-7_054-v0010}{10}
\psnumber{ps_134-7_054-v0011}{11}
\psnumber{ps_134-7_054-v0012}{12}
\psnumber{ps_134-7_054-v0013}{13}
% Pursuing Stacks (cotes 134-2 to 134-8): archive cote 134-7, site page 54 (PDF page 55), typescript page 496.
% typescript. Draft compared: ps6.tex:2663-2741.
% AI-transcribed by Claude Opus 5.5 (claude-opus-5-5, adaptive thinking, effort high), Message Batches API, batch ps4 (msgbatch_012Aoc1NuskATPsyBeKSYK7J), 2026-11-12; reviewed against both supplied scan halves; see CHECK_DECISIONS.tsv.
\scanpage{134-7/54}
\mspage{496}

basic structures as $K(\underline{Z},n)$ with $n\hand{\geqslant} 2$. At any rate, the \hand{``}way-out\hand{''} argument
I have finally been thinking of, in order to try and check (6) is a quasi%
isomorphism, rests on vanishing assumptions which (as I was informed by
Illusie the next day) are wholly unrealistic. This finally clears up,
it would seem, a tenacious misconception which has been sticking to my
first heuristic ponderings about the homology and cohomology formalism
for schematic homotopy types: \uline{one should be very careful not to substitu\var{\xout{t}}{}{typewriter deletion}{ps_134-7_054-v0001}}%
\uline{te the ``pointed'' linearization} $L_{pt}(W(M))\hand{\simeq}\hand{\Gamma}_k(M)$ \uline{for the nonpointed}
\uline{one} $L(W(M))\hand{\simeq}\hand{\Gamma}_k{}^{\hand{\wedge}}(M)$, \uline{in computing homology and cohomology invariants}
\uline{of schematic homotopy types}. To say it differently, in order to be
able to \var{\xout{define}}{}{typewriter deletion}{ps_134-7_054-v0002} compute (or just define) the ``schematic'' homology and
cohomology invariants, we do need as a model a fullfledged semisimplicial
unipotent bundle, not just the corresponding formal one, giving rise
to invariants of \var{\deleted{their}}{}{deletion}{ps_134-7_054-v0003}\var{\inserted{it's}}{\hand{it's}}{insertion}{ps_134-7_054-v0004} own, namely ``formal'' homology and cohomology,
which are definitely distinct from the former.

\var{\deleted{\xout{It is all the more remarkable that}}}{}{deletion}{ps_134-7_054-v0005}
118. It is all the more remarkable, in view of the preceding findings,
that the homotopy invariants
\[ \hand{\pi}_i(X_*) \ , \quad i \hand{\geqslant} 0 , \]
of a \var{\inserted[typed]{pointed}}{pointed}{typed insertion}{ps_134-7_054-v0006} semisimplicial \var{\inserted[typed]{unipotent}}{unipotent}{typed insertion}{ps_134-7_054-v0007} bundle $X_*$ (still assuming the components $X_n$ to be
smooth, i.e. \var{\inserted{to}}{\hand{to}}{insertion}{ps_134-7_054-v0008} correspond\var{\deleted{ing}}{}{deletion}{ps_134-7_054-v0009} to projective modules) turn out to be invari-%
ants of the corresponding ``formal'' object, and, more startling still,
of the corresponding ``infinitesimal'' object of order 1. M\var{\over[typed]{g}{o}}{o}{typewriter overwriting}{ps_134-7_054-v0010}re specifical-%
ly, consider the ``Lie functor'' or ``tangent space at the origin''
\[ (7) \qquad \mathrm{Lie}: U(k)^{\hand{\bullet}} \hand{\longrightarrow} Ab_k \ , \qquad X = W(M) \hand{\longmapsto} M \ , \]
which we'll need only for the time being for smooth $X$, when the geometri
meaning of it is clear. This functor transforms semisimplicial \var{\inserted[typed]{pointed}}{pointed}{typed insertion}{ps_134-7_054-v0011} bundles
into semisimplicial $k$-modules $\mathrm{Lie}(X_*)$ , thus we should get, besides
abelianization, another remarkable functor\hand{,} from $\hand{\mathcal{M}_1}(k)$ \var{\inserted{say}}{\hand{say}}{insertion}{ps_134-7_054-v0012} to $D.(Ab_k)$ :
\[ (8) \qquad \mathrm{Lie}: \hand{\mathcal{M}_1}(k) \hand{\longrightarrow} D.(Ab_k) \ \var{\inserted[typed]{\text{via}}}{\text{via}}{typed insertion}{ps_134-7_054-v0013}, \ X_* \hand{\longmapsto} \mathrm{Lie}(X_*) \ , \]
granting that $X_* \hand{\mapsto} \mathrm{Lie}(X_*)$ transforms quasiisomorphisms into quasiiso-%
morphisms. Now that this must be so follows from the really startling
formula
\[ (9) \qquad \hand{\pi}_i(X_*) \hand{\simeq} \hand{\pi}_i(\mathrm{Lie}(X_*)) \quad (\hand{\simeq} H_i \text{ of the associated chain} \]
\[ \text{complex in } Ab_k \ ) \ , \]
where the left hand side, I recall, is defined as
\[ \hand{(10)} \qquad \hand{\pi}_i(X_*) = \hand{\pi}_i(X_*(k)) \ . \]

\margin{54}

% Visual collation: both supplied image halves inspected; see CHECK_DECISIONS_134_7_8.tsv.

\psapparatusstart
\psentry{ps_134-7_054-v0001}{1}{typewriter deletion}{\xout{t}}{}{The typewriter strikes out this material; it leaves the corrected running text.}
\psentry{ps_134-7_054-v0002}{2}{typewriter deletion}{\xout{define}}{}{The typewriter strikes out this material; it leaves the corrected running text.}
\psentry{ps_134-7_054-v0003}{3}{deletion}{\deleted{their}}{}{The author cancels this material by hand; it leaves the corrected running text.}
\psentry{ps_134-7_054-v0004}{4}{insertion}{\inserted{it's}}{\hand{it's}}{The handwritten insertion is incorporated at its indicated position; insertion brackets leave the running text.}
\psentry{ps_134-7_054-v0005}{5}{deletion}{\deleted{\xout{It is all the more remarkable that}}}{}{The author cancels this material by hand; it leaves the corrected running text.}
\psentry{ps_134-7_054-v0006}{6}{typed insertion}{\inserted[typed]{pointed}}{pointed}{The typed insertion is incorporated at its indicated position; insertion brackets leave the running text.}
\psentry{ps_134-7_054-v0007}{7}{typed insertion}{\inserted[typed]{unipotent}}{unipotent}{The typed insertion is incorporated at its indicated position; insertion brackets leave the running text.}
\psentry{ps_134-7_054-v0008}{8}{insertion}{\inserted{to}}{\hand{to}}{The handwritten insertion is incorporated at its indicated position; insertion brackets leave the running text.}
\psentry{ps_134-7_054-v0009}{9}{deletion}{\deleted{ing}}{}{The author cancels this material by hand; it leaves the corrected running text.}
\psentry{ps_134-7_054-v0010}{10}{typewriter overwriting}{\over[typed]{g}{o}}{o}{The superposed typed reading replaces the underlying reading.}
\psentry{ps_134-7_054-v0011}{11}{typed insertion}{\inserted[typed]{pointed}}{pointed}{The typed insertion is incorporated at its indicated position; insertion brackets leave the running text.}
\psentry{ps_134-7_054-v0012}{12}{insertion}{\inserted{say}}{\hand{say}}{The handwritten insertion is incorporated at its indicated position; insertion brackets leave the running text.}
\psentry{ps_134-7_054-v0013}{13}{typed insertion}{$\inserted[typed]{\text{via}}$}{$\text{via}$}{The typed insertion is incorporated at its indicated position; insertion brackets leave the running text.}
\psapparatusend

\end{document}
